    \section{The Algorithm}

    We now summarize the algorithm in two parts. The main procedure first performs the preprocessing phase, in which it computes the partially rounded \textsc{LP} solution $(\hat{x},\hat{y})$ and the vector solution $\hat{v}$ obtained from the \textsc{SDP}, and then it invokes the recursive routine. The recursive routine only builds the clustering tree, using the guides prepared by the main procedure.

    \begin{theorem}\label{thm:main-approximation}
        If $\cF$ is a $\alpha\br{n}$-well-behaved class of graphs, then there exists an $O\br{\alpha\br{n}+\beta\br{n}}$-approximation algorithm for the \textsc{HC-$\cF$} problem.
    \end{theorem}
    \begin{proof}

    
    Denote by $\beta\br{n}$ the approximation ratio of the rounding algorithm for \textsc{$k$-BP}. 

    \begin{algorithm}[H]
\caption{Main \textsc{HC-$\cF$} Procedure}
\label{alg:main-procedure}
\label{alg:main-hac}
\DontPrintSemicolon
\KwIn{A weighted graph $G=(V,E,c,w)$}
\KwOut{A clustering tree for $G$}
\tcp{Preprocessing}
Solve \textsc{LP-HC-$\cF$} on $G$ and obtain a fractional solution $(\mathbf{x},\mathbf{y})$\;
\ForEach{$t\in [w\br{G}]$ and $i\in V$}{
    Set $\hat{y}_i^t \gets 1$ if $y_i^t \geq 1/2$, and set $\hat{y}_i^t \gets 0$ otherwise\;
}
\ForEach{$t\in [w\br{G}]$ and $uv\in E$}{
    $\hat{x}_{uv}^t\gets 2\cdot x_{uv}^t$\;
}
\tcp{Clustering phase}
\Return{\textsc{Recursive-HC-$\cF$}($G,\hat{\mathbf{x}},\hat{\mathbf{y}})$}\;
\end{algorithm}

    \begin{algorithm}[H]
\caption{Recursive \textsc{HC-$\cF$} Procedure}
\label{alg:recursive-procedure}
\label{alg:recursive-hac}
\DontPrintSemicolon
\KwIn{A connected subgraph $H$, together with the guides $(\hat{\mathbf{x}},\hat{\mathbf{y}})$ restricted to $V(H)$}
\KwOut{A clustering tree for $H$}
\If{$H\in \cF$}{
    \Return{a clustering tree $T$ consisting of a single node labeled by $H$}\;
}
Set $r\gets w\br{V\br{H}}$ and $t\gets r/2$\;
Set $X_1\gets \brc{i\in V\br{H}: \hat{y}_i^t = 1}$\;
Run the \textsc{$p_{\cF}$-P} rounding algorithm on \textsc{LP-$p_{\cF}$-P} guide on $H$ at level $t$ from $(\hat{\mathbf{x}},\hat{\mathbf{y}})$ with $X=X_1$ and let $F_H\subseteq E(H)$ be the returned cut\;

Set $X_0\gets \brc{i\in V\br{H}: \hat{y}_i^t = 0}$\;
Run the weighted \textsc{$\rho$-SEP} rounding algorithm with $\rho=1/2$ and $\rho_0=2/3$ on \textsc{LP-$\rho$-SEP} guide on $H$ at level $t$ from $\hat{\mathbf{x}}$ with $X=X_0$ and let $B_H\subseteq E(H)$ be the returned cut\;

$S_H\gets F_H\cup B_H$\;

Initialize a root node $u_H$ of the output tree and label it by $S_H$\;
\ForEach{$A\in H\setminus S_H$}{
    Recursively compute $T_A\gets $\textsc{Recursive-HC-$\cF$}($A,\hat{\mathbf{x}},\hat{\mathbf{y}})$\;
        Attach $T_A$ as a child of $u_H$\;
}
\Return{the clustering tree rooted at $u_H$}\;
\end{algorithm}

    \begin{lemma}\label{lem:recursive-cost}
        Fix some cluster $H$ for which the recursive procedure was called. Then, 
        \[
        \COST_G\br{T, u_H}\leq 8\alpha\br{n}\cdot \sum_{t=\spr{H}/8+1}^{\spr{H}} \COST\br{LP_t^H}+8\beta\br{n}\cdot \sum_{t=\spr{H}/8+1}^{\spr{H}} \COST\br{SDP_t^H}.
        \]
    \end{lemma}
    \begin{proof}
        Firstly, consider the contribution of edges in $F_H$ to the cost. Since it it being cut in a cluster of size $\spr{H}$, we have that it is bounded by $\spr{H}\cdot c\br{F_H}$. By applying Lemma \ref{lem:lp-sdp-bound} with $d=8$ we get that:
        \begin{align*}
            \spr{H}\cdot c\br{F_H} \leq \spr{H}\cdot \COST\br{LP_{\spr{H}/4}^H} \leq 8\alpha\br{n}\cdot \sum_{t=\spr{H}/8+1}^{\spr{H}/4} \COST\br{LP_t^H}.
        \end{align*}
        Secondly consider the contribution of all $B_H$ to the cost. Since it it being cut in a cluster of size $\spr{H}$, we have that it is bounded by $\spr{H}\cdot c\br{B_H}$. Applying Lemma \ref{lem:lp-sdp-bound} again with $d=8$ we get:
        \begin{align*}
             \spr{H}\cdot c\br{B_H} \leq \spr{H}\cdot \COST\br{SDP_{\spr{H}/4}^H} 
             \leq 8\beta\br{n}\cdot \sum_{t=\spr{H}/8+1}^{\spr{H}/4} \COST\br{SDP_t^H}
        \end{align*}
    Since we know that $c\br{S_H}= c\br{F_H\cup B_H}\leq c\br{F_H}+c\br{B_H}$, and we have $\COST_G\br{T, u_H} \leq \spr{H}\cdot c\br{S_H}$, the claim follows.
    \end{proof}
    Armed with the above lemma we are ready to sum the contribution of all clusters to the cost of the solution. Let $T$ be the clustering tree returned by the algorithm and let $\cC_T$ be the set of all non-tree clusters for which the recursive procedure was called. We have:
    \begin{align*}
        \COST_G\br{T} &= \sum_{H \in \cC_T}\COST_G\br{T, u_H}
        \\&\leq 8\cdot\sum_{H \in \cH} \sum_{t=\spr{H}/8+1}^{\spr{H}/4} \br{\alpha\br{n}\cdot \COST\br{LP_t^H} + \beta\br{n}\cdot \COST\br{SDP_t^H}}\\
        &\leq 8\cdot\br{\alpha\br{n}\cdot\COST\br{LP} + \beta\br{n}\cdot\COST\br{SDP}}\\
        &\leq 16\cdot\br{\alpha\br{n}+\beta\br{n}}\cdot \OPT
    \end{align*}
     where the first inequality is due to Lemma \ref{lem:recursive-cost}, the second inequality is due to Lemma \ref{lem:relaxation-decomposition} and the third inequality is due to Lemma \ref{lem:lp-rounding}. 
    \end{proof}

We get the following corollaries:
\begin{corollary}
    There exists an $O\br{\log n \cdot \log\log n}$-approximation algorithm for the \textsc{HC-$\cT$} problem.
\end{corollary}
\begin{corollary}
    There exists an $O\br{\log n}$-approximation algorithm for the \textsc{HC-$\cD_d$} problem.
\end{corollary}
